%==============================================================================
%  sections_v3/proofs_existence.tex -- the paper proof of Proposition P1
%  (existence of a cutoff perfect Bayesian equilibrium).
%
%  SOURCE OF TRUTH for the STATEMENT: research/model_v4/MODEL_CARD.md, version
%  stamp 2026-08-30 (F5 route R40-A, gate PASS, theory-record freeze;
%  commit 65b8db3), section 6 P1 row, as
%  transcribed into sections_v3/theorem_section.tex.  The mathematics is
%  transcribed from research/model_v4/proofs/P1_proof.md (repairs applied
%  through ticket 35 rounds 1-2 and the close-out sweep), whose hypothesis
%  numbering h.1-h.17 maps onto the statement's (P-1)-(P-13) as follows:
%
%    h.1  = A1                 -> (P-1)        h.10 = no-feedback + flag term.  -> (P-8)
%    h.2  = A2'                -> (P-2)        h.11 = round-2 action set        -> (P-9)
%    h.3  = A3                 -> (P-3)        h.12 = m_0 >= 0                  -> (P-11)
%    h.4  = A4                 -> (P-4)        h.13 = H-ord      -> final A8 paragraph
%    h.5  = STRUCK (A5)  -> "A5 is not assumed" trailing paragraph
%    h.6  = A6                 -> (P-5)        h.14 = U_j objective             -> (P-12)
%    h.7  = A7-J               -> (P-6)        h.15 = upper-set engagement flag -> final A8 par.
%    h.8  = A8           -> final A8 paragraph h.16 = continuation-cost equiv.  -> (P-10)
%    h.9  = D1 by statement    -> (P-7)        h.17 = TR block                  -> (P-13)
%
%  Step numbering is 1:1 with the source (Steps 1-20, Parts A-F); nothing is
%  merged or renumbered, so every audit citation into the source resolves here.
%  The source's structural meta-blocks (hypothesis register, NOT CLAIMED list,
%  NOTATION DELTA, repair tags) are process and do not appear; their substance
%  travels as in-proof sentences and as the closing remark paragraph.
%
%  Internal labels are all prefixed eq:P1- .  No \section and no sec:proofs
%  label lives here: sections_v3/proofs_section.tex owns both.
%==============================================================================

\begin{proof}[Proof of Proposition~\ref{prop:P1} (P1: existence)]
The argument builds the equilibrium object rather than characterising it. Parts~A and~B fix a
conjectured cutoff vector and solve the price system it induces, separating a finite pooled layer
from a continuum-indexed flagged layer; Part~C supplies beliefs at every history that carries a
requirement, on path and off; Part~D discharges sequential optimality of the flagged component at
every flagged pair; Part~E builds the outer best-response map, applies Brouwer's fixed-point theorem
and assembles the six requirements of Definition~\ref{def:cpbe}; Part~F is the addendum under
Assumption~\ref{asm:A8} (A8). The parameter vector $\vartheta$ is fixed throughout and suppressed where it
does not vary.

\smallskip\noindent\emph{Notation used only in this proof.} Write $\sigma_F$ for a generic value of
the flagged tuple $\Sfl$ of Definition~\ref{def:prices}, and
$[\underline s,\overline s]$ for the common signal bracket underlying the polytope $\Theta$ of
Definition~\ref{def:theta}. Write $\Phi_s$ and $\varphi_s$ for the c.d.f.\ and density of the signal
$s$, reserving $\Phi$ for the standard normal c.d.f.\ of \eqref{eq:entry} and $\phi$ for its density.
At a control-node information set $\mathcal I$ put
\[
\hat v(\mathcal I) = \Eop[v \mid \mathcal I],
\qquad
\pi(\mathcal I) = \Prb(a = 1 \mid \mathcal I),
\qquad
\bar m(\mathcal I) = m_0 + \pi(\mathcal I)\,\Delta_m .
\]
Further symbols --- $\mathcal P_{\mathcal I}$, $\varrho$, $A$, $\mathcal G_F$, $\iota_F$, $j_k$,
$j^\star$, $\mathcal S(k)$, $\mathfrak w$, $\mathcal Q_j(s)$, $G_j$, $E_j$, $\mu_n$, $L_j$, $w_n$,
$t_n$, $Z_n$, $\Lambda_k$, $\Lambda_u$, $(\hat v_\circ,\pi_\circ)$, $P_\circ$,
$\operatorname{supp}(z_d)$, $s_F$ --- are declared where they are first used. The engagement cost
$C_j(s)$ and the execution outlay $\mathcal C_j^{\mathrm{trade}}$ of Definition~\ref{def:U} are
always written with their subscripts, so that neither collides with the chord $C_h$ of
\eqref{eq:chord}, the composition ratios $C_\tau,C_T$, or the cells $\mathcal C_F,\mathcal C_P$.

\medskip
\noindent\emph{Part A: the game at a fixed conjecture.}

\smallskip
\noindent\textbf{Step 1 (the conjecture induces a measurable plan-selection map).}
Fix a conjectured cutoff vector $k = (k_1 \le \dots \le k_{J-1})$ in
\[
\Theta \;=\; \bigl\{k \in [\underline s,\overline s]^{J-1} :
\underline s \le k_1 \le \dots \le k_{J-1} \le \overline s\bigr\},
\]
the cutoff polytope of Definition~\ref{def:theta}. It is nonempty, compact and convex, being the
intersection of a cube with the $J-2$ half-spaces $\{k_i \le k_{i+1}\}$; nonemptiness is not
decoration, since Brouwer's theorem is vacuous without it, and it is supplied by
Definition~\ref{def:theta}. Define the map induced by the conjecture,
\begin{equation}
\label{eq:P1-jk}
j_k(s) \;=\; 1 + \#\bigl\{i \in \{1,\dots,J-1\} : k_i \le s\bigr\}.
\end{equation}
Each set $\{s : k_i \le s\}$ is a half-line, so $j_k$ is a weakly increasing step function of $s$
with values in $\mathcal J$, and it is Borel. This is the object requirement~(i) of
Definition~\ref{def:cpbe} calls a weakly ordered cutoff vector mapping the signal into a plan. The
second clause of Assumption~\ref{asm:A3} (A3) --- hypothesis (P-3), that the preferred plan is weakly
increasing in $s$ --- is what makes a representation of this shape the right one for a best response;
Step~13 returns to that point and derives the representation rather than assuming it.

\smallskip
\noindent\textbf{Step 2 (every date-$0$ object is a deterministic Borel function of the pair
$(j,s)$).}
By hypothesis (P-8), the no-feedback timing of Assumption~\ref{asm:nofeedback}, the executed path
carries no feedback from realised order flow or from realised prices. Hence for each $j \in
\mathcal J$ the objects $B_j(s,d)$ for $d = 0,\dots,H$, the pooled marks $q_{jd}(s) =
\Gamma\bigl(B_j(s,d)-B_j(s,d-1)\bigr)$, the crossing date $c_j(s;\tau)$, the filing date $f_j(s) =
c_j(s)+T$, the stake at filing $B_j^F(s) = B_j(s,f_j(s))$ and the flagged order $Q_j^F(s) =
b_j^*(s)-B_j^F(s)$ are functions of $(j,s)$ and of the policy pair $(\tau,T)$ alone.

Measurability in the signal is where care is needed, and it is consumed from the primitive table
rather than inherited. That $s \mapsto B_j(s,d)$ is Borel \emph{for every plan, Exit included}, is
clause (TR-ii) of Assumption~\ref{asm:TR} (TR), carried by hypothesis (P-13) and consumed here
\emph{directly}: the monotonicity $\partial_s B_j \ge 0$ of clause (TR-iii) is imposed on Voice
alone, Exit paths are restricted in $d$ and unrestricted in $s$, and the conclusion of
Lemma~\ref{lem:D1} is measurability of $D$ and of the cell map, not of the objects listed above. The
coarsening $\Gamma$ has finite image by Assumption~\ref{asm:A2p} (A2$'$), hypothesis (P-2), hence is Borel.
The crossing date $c_j(\cdot;\tau) = \inf\{d : B_j(\cdot,d) \ge \tau\}$ is the pointwise minimum,
over the finite calendar supplied by the same assumption, of the indices of the Borel sets
$\{B_j(\cdot,d) \ge \tau\}$, so it is Borel with values in $\{0,\dots,H\} \cup \{+\infty\}$. The
stake at filing is a finite sum of products of Borel functions,
\begin{equation}
\label{eq:P1-BF}
B_j^F(s) \;=\; \sum_{d=0}^{H-T} \ind\{f_j(s) = d+T\}\; B_j(s,d),
\end{equation}
hence Borel, and so is $Q_j^F$. By hypothesis (P-7), Lemma~\ref{lem:D1} with its own hypotheses
travelling, the disclosure indicator is $D_j(s;\tau,T) = \ind\{a_j = 1\}\,\ind\{B_j(s,H-T) \ge
\tau\}$ --- the clock equivalence of part~(b) of that lemma --- and is Borel in $s$. Composing with
\eqref{eq:P1-jk}, all of these become Borel functions of the signal alone at the fixed conjecture
$k$. The definitions of $c_j$, $f_j$, $B_j^F$, $Q_j^F$ and $b_j^*$ used here are clause (TR-iii)'s,
again by hypothesis (P-13).

\smallskip
\noindent\textbf{Step 3 (the pooled public-history family is finite; the flagged family is not).}
Definition~\ref{def:prices} sets $\mathcal H_d^P = (X_0,\dots,X_d;\ \text{flag landed by } d)$ with
$X_d = q_{jd}+z_d$. By hypothesis (P-2) the image of $\Gamma$ is finite and $z_d$ takes values in
$\{-\bar z,0,+\bar z\}$, so each $X_d$ takes values in a finite set; the calendar
$\{0,\dots,H\}$ is finite; and the flag coordinate is a single bit. The collection of pooled public
histories is therefore finite.

The flagged family is not. By Step~2 the flagged tuple $\sigma_F$ is Borel, but it takes values in
$[0,\bar b]^2 \times \{1\}$, a continuum: Definition~\ref{def:plans} puts $B_j(s,d) \in [0,\bar b]$
with $s$ Gaussian, clause (TR-ii) records that stake levels are continuum-valued, and the finiteness
of Assumption~\ref{asm:A2p} covers the plan menu, the image of $\Gamma$, the noise support and the
calendar --- not the stake level. The two layers must therefore be handled differently, and Steps~5
and~6 do so. The finiteness just established rests on Definition~\ref{def:prices}'s own content, in
which the flag enters $\mathcal H_d^P$ as a bit and the filing content $B^F$ does not; were
post-filing pooled histories to carry $B^F$, the pooled family would join the continuum and the
selection argument of Step~6 would have to be run there as well.

\medskip
\noindent\emph{Part B: inner prices.}

\smallskip
\noindent\textbf{Step 4 (every control-node pricing fixed point reduces to one scalar equation, and
depends on the information set only through the pair $(\hat v,\pi)$).}
Fix a control-node information set $\mathcal I$. By hypothesis (P-13), clause (TR-iv) of
Assumption~\ref{asm:TR}, the terminal payoff is \eqref{eq:Y}, prices obey $P(\mathcal I) =
\Eop[Y \mid \mathcal I]$ and entry obeys \eqref{eq:entry}; writing $\mathsf B$ for the entry
indicator, \eqref{eq:entry} is $\mathsf B = \ind\{\xi \ge P(\mathcal I)+K+\bar m(\mathcal I)-\bar
S\}$. Given $\mathcal I$, the quantities $P(\mathcal I)$ and $\bar m(\mathcal I)$ are
$\mathcal I$-measurable, so $\mathsf B$ is a function of $\xi$ alone. By hypothesis (P-1),
Assumption~\ref{asm:A1} (A1), $\xi$ is independent of $(v,\varepsilon)$ and of every $z_d$, hence of
$(v,s,z_{0:H})$ jointly and therefore of $(v,a,\mathcal I)$; conditionally on $\mathcal I$ it is
independent of $(v,a)$. Writing $p = \Prb(\mathsf B = 1 \mid \mathcal I)$ and taking conditional
expectations of \eqref{eq:Y} term by term,
\begin{equation}
\label{eq:P1-EY}
\Eop[Y \mid \mathcal I]
\;=\; (1-p)\bigl(\hat v + \pi\Delta_V\bigr) + p\,(P+m_0) + \Delta_m\,p\,\pi
\;=\; (1-p)\bigl(\hat v + \pi\Delta_V\bigr) + p\bigl(P + \bar m\bigr).
\end{equation}
With $\xi \sim N(0,\sigma_\xi^2)$, which is clause (TR-i)'s distributional form under the same
hypothesis, $p = 1-\Phi\bigl((P+K+\bar m-\bar S)/\sigma_\xi\bigr)$, which is \eqref{eq:entry}
verbatim and lies in $(0,1)$ for every finite $P$. Define the inner pricing map
\begin{equation}
\label{eq:P1-inner}
\mathcal P_{\mathcal I}(P) \;=\; \bigl(1-p(P)\bigr)\bigl(\hat v + \pi\Delta_V\bigr)
+ p(P)\bigl(P+\bar m\bigr),
\qquad
p(P) = 1-\Phi\!\Bigl(\tfrac{P+K+\bar m-\bar S}{\sigma_\xi}\Bigr).
\end{equation}
The requirement $P(\mathcal I) = \Eop[Y \mid \mathcal I]$ is then the scalar equation
$\mathcal P_{\mathcal I}(P) = P$, and the map depends on $\mathcal I$ only through the two scalars
$\bigl(\hat v(\mathcal I),\pi(\mathcal I)\bigr)$. That reduction is what Steps~5 to~8 run on, and it
is the object the trailing paragraph of the proposition calls the scalar reduction.

\smallskip
\noindent\textbf{Step 5 (the pooled layer, in two parts, because only one of them is a fixed point).}
By Step~3 there are finitely many pooled public histories. Step~4's map was derived at a control
node, which is where $\mathsf B$ is a function of $\xi$ alone given the conditioning, so the two
pooled layers are treated separately.

\emph{(a) The pooled control-node cell, $D = 0$ at date $H$.} Here $\mathcal I = \mathcal I_H$ is a
control node and Step~4 applies as derived. Step~7 supplies a unique fixed point of
$\mathcal P_{\mathcal I}$ from hypothesis (P-11), the restriction $m_0 \ge 0$ of \eqref{eq:m0}, and
Step~8 supplies its continuity in the belief. This is a genuine fixed point: the price appears on
both sides of \eqref{eq:P1-inner} through the entry indicator. Neither clause is taken from
Assumption~\ref{asm:A5} (A5), which the proposition does not assume; Steps~7 and~8 do not depend on the
present step, so citing them forward is not circular.

\emph{(b) Intermediate pooled dates $d < H$.} A history $\mathcal H_d^P$ with $d < H$ is not a
control node. Display \eqref{eq:Y} writes the takeover branch as $\mathsf B(P+m_0+a\Delta_m)$ with
$P$ unqualified. Under the reading Step~4 itself adopts, that $P$ is the control-node price
$P(\mathcal I_H)$, so by the tower property
$P_d^P = \Eop\bigl[Y \mid \mathcal H_d^P\bigr]$ is a plain conditional expectation of
already-solved control-node values, with no self-reference and no fixed point --- which is the
reading the paragraph after Definition~\ref{def:prices} records. Under the alternative reading, in
which the $P$ inside $Y$ is the price at whichever information set is conditioning, part~(a)'s
fixed-point argument applies at these dates too. The conclusion this step is used for survives on
either reading, and that is why nothing downstream turns on the adjudication.

That conclusion is: a finite family requires no selection argument, and the pooled price family
$k \mapsto \bigl(P_d^P(\mathcal H_d^P;k)\bigr)_{\mathcal H_d^P}$ is a finite vector of continuous
functions of $k$ --- at the control-node cell by (a) with Steps~7 and~8, and at $d < H$ by (b) as a
finite-sum conditional expectation of continuous functions --- on those histories carrying positive
probability under the conjecture $k$. Histories null under $k$ but reachable, and histories null
under every profile, are supplied in Step~9(b) and Step~9(c) respectively.

\smallskip
\noindent\textbf{Step 6 (the flagged layer is a measurably selected family, built rather than
assumed).}
By Step~3 the flagged information sets are indexed by the continuum $\sigma_F \in [0,\bar b]^2 \times
\{1\}$. A pointwise statement --- at each $\sigma_F$ there is a unique root --- does not by itself
yield a \emph{function} of $\sigma_F$ that the model can integrate against, which is what
$M_F = \Delta_m\Eop[h \mid D=1]$ of Definition~\ref{def:premium} and the timing split of
Lemma~\ref{lem:D1} both require. The family is constructed in four moves.

\emph{(a) The engagement posterior is degenerate on the flagged cell.} Assumption~\ref{asm:A4} (A4),
hypothesis (P-4), together with clause (TR-iii) gives $D = 1 \Rightarrow a = 1$, and hypothesis (P-7)
makes $\{D = 1\}$ an event of the control-node history, so $\Prb(a = 1 \mid \sigma_F, D=1) = 1$. This
is the row of Definition~\ref{def:prices} that sets $\pi = 1$ on $\mathcal C_F$, and hence
$\bar m = m_0+\Delta_m = m_1$ on the whole flagged cell, a constant.

\emph{(b) The flagged pricing map therefore depends on one scalar.} By Step~4 and (a), the flagged
map depends on $\sigma_F$ only through $\hat v(\sigma_F;k) = \Eop[v \mid \sigma_F, D=1]$. Write
$\mathcal G_F$ for the map sending a belief $\hat v$ to the root of
$\mathcal P_{\mathcal I}(\cdot)-\mathrm{id}$ at $(\hat v,\pi=1)$. Step~7 makes $\mathcal G_F$
single-valued, by existence and uniqueness of the root under hypothesis (P-11), and Step~8 makes it
continuous in the belief --- indeed $1$-Lipschitz, since $\partial P/\partial\hat v \in (0,1]$ there.
This is the one place at which continuity in the belief is load-bearing, because (d) below composes a
Borel map with a continuous one.

\emph{(c) The belief is a Borel function of the tuple.} Two routes are available and both are taken.
First, $\hat v(\cdot;k)$ is a conditional expectation with respect to $\sigma(\sigma_F)$ and is
therefore $\sigma(\sigma_F)$-measurable by construction. Second, under hypothesis (P-6),
Assumption~\ref{asm:A7J} (A7-J), joint tuple injectivity, the map $(j,s) \mapsto \sigma_F$ is
injective on the flagged-pair set and, by Step~2, Borel; both $\mathcal J \times \mathbb R$ and
$[0,\bar b]^2\times\{1\}$ are Borel subsets of Polish spaces, so the Lusin--Souslin theorem supplies
a Borel inverse $\iota_F$ on the image, and
$\hat v(\sigma_F;k) = \mu_v + \beta\bigl(\iota_F(\sigma_F)_s - \mu_v\bigr)$
with $\beta$ the Gaussian projection coefficient of Definition~\ref{def:primitives}. Injectivity
together with the Borel regularity of clause (TR-ii) already delivers the measurable inverse; no
separate selection assumption is introduced, which is the sense in which the measurable-selection
content of Assumption~\ref{asm:A5} is derived here from Assumption~\ref{asm:A7J} (A7-J) and clause (TR-ii)
rather than assumed.

\emph{(d) The family.} Therefore $P^F(\sigma_F;k) = \mathcal G_F\bigl(\hat v(\sigma_F;k)\bigr)$ is
the composition of a Borel map with a continuous map, hence Borel. This is the measurably selected
flagged price family, and it is pinned rather than chosen: uniqueness of the root at each $\sigma_F$
leaves no freedom, so no selection principle is invoked and no two runs of the argument can produce
different families. That is a statement about the \emph{price} family given the belief; the sense in
which the flagged \emph{belief} is pinned is Step~10's, and is weaker.

\smallskip
\noindent\textbf{Step 7 (under $m_0 \ge 0$ the inner root exists and is unique, by derivation).}
Write $A = \hat v + \pi\Delta_V$ for the no-takeover branch value at $\mathcal I$ --- a symbol used
only in this step and the next, and never subscripted, so that it does not collide with the weights
$A_0,A_{1/2},A_1$ of \eqref{eq:atau} --- and put
\begin{equation}
\label{eq:P1-rho}
\varrho(P) \;=\; \mathcal P_{\mathcal I}(P) - P \;=\; \bigl(1-p(P)\bigr)(A-P) + p(P)\,\bar m,
\end{equation}
continuous in $P$ because $\Phi$ is. By hypothesis (P-11), the restriction $m_0 \ge 0$ of
\eqref{eq:m0}, together with $\Delta_m > 0$ and $\pi \in [0,1]$ from clause (TR-i), we have
$\bar m \ge 0$.

\emph{(i) No root below $A$.} For $P < A$ both terms of \eqref{eq:P1-rho} are nonnegative and the
first is strictly positive, since $p(P) < 1$ by Step~4; hence $\varrho(P) > 0$.

\emph{(ii) A root exists.} We have $\varrho(A) = p(A)\bar m \ge 0$. If $\bar m = 0$ then $P = A$ is a
root. If $\bar m > 0$ then $\varrho(A) > 0$, while as $P \to +\infty$ we have $p(P) \to 0$ and
$(A-P) \to -\infty$, so $\varrho(P) \to -\infty$. An explicit bracket: for
$P \ge \bar S-K-\bar m+\sigma_\xi$ one has $p(P) \le 1-\Phi(1) < 0.159$, whence
$\varrho(P) \le 0.159\,\bar m - 0.841\,(P-A) \le 0$ once additionally $P \ge A+0.19\,\bar m$. The
intermediate value theorem on $\bigl[A,\ \max\{\bar S-K-\bar m+\sigma_\xi,\ A+0.19\bar m\}\bigr]$
delivers a root.

\emph{(iii) The root is unique.} The residual is differentiable, with
$\varrho'(P) = p'(P)\bigl(P+\bar m-A\bigr) + p(P) - 1$ and
$p'(P) = -\phi\bigl((P+K+\bar m-\bar S)/\sigma_\xi\bigr)/\sigma_\xi < 0$. At any root, (i) gives
$P \ge A$, so $P+\bar m-A \ge \bar m \ge 0$ and the first term is at most zero, while the second is
strictly negative because $p < 1$. Hence $\varrho' < 0$ \emph{strictly at every root}. Suppose there
were two roots $P_1 < P_2$ with no root between them. Then $\varrho'(P_1) < 0$ forces $\varrho < 0$
immediately to the right of $P_1$, and since $\varrho$ has no zero on $(P_1,P_2)$ it is negative
throughout that interval; but $\varrho'(P_2) < 0$ forces $\varrho > 0$ immediately to the left of
$P_2$. The two conclusions contradict each other, so there is at most one root.

On the maintained sign $m_0 \ge 0$ the existence and uniqueness content of Assumption~\ref{asm:A5} is
therefore a theorem rather than an assumption, which is the first of the three parts of the
proposition's trailing paragraph on that assumption. What is left over is continuity of the
\emph{composition} in the conjecture $k$, which runs through the conditioning pair $(\hat v,\pi)$
rather than through the pricing map; Step~15 takes that up and says exactly where it is assumed.

\smallskip
\noindent\textbf{Step 8 (the inner root is monotone and non-expansive in the belief).}
Continuity in the belief comes from the same scalar reduction. On \eqref{eq:P1-rho} the residual is
continuously differentiable jointly in $(P,\hat v)$, and by Step~7(iii) $\partial_P\varrho < 0$
strictly at every root, so the implicit-function argument applies to $\varrho(P;\hat v) = 0$ and
yields
\begin{equation}
\label{eq:P1-dP}
\frac{\partial P}{\partial \hat v}
\;=\; \frac{1-p}{\,1-p+\lvert p'(P)\rvert\,(P+\bar m-A)\,}
\;\in\; (0,1],
\end{equation}
the denominator being at least $1-p > 0$ by hypothesis (P-11) and Step~7(i). So the root map is
single-valued, strictly increasing and $1$-Lipschitz in the belief. This is the continuity content of
Assumption~\ref{asm:A5}, derived rather than assumed, and it is what Step~6(b) consumes.

\medskip
\noindent\emph{Part C: beliefs, on path and off.}

\smallskip
\noindent\textbf{Step 9 (pooled beliefs as limits of one full-support perturbation family, on the
reachable history set, at every $\kappa \in [0,1]$).}
Requirement~(vi) of Definition~\ref{def:cpbe} asks for off-path beliefs as limits of full-support
perturbations. The family is fixed once here and used at every $k \in \Theta$, not only at the fixed
point of Step~16 --- which is what the deviation payoffs defining the outer map require, since those
payoffs read off-path pooled histories.

\emph{(a) The alphabet and the reachable set.} Definition~\ref{def:primitives} sets
$\Prb(z_d = 0) = 1-\kappa$ and $\Prb(z_d = \pm\bar z) = \kappa/2$ on the whole domain
$\kappa \in [0,1]$, so the \emph{realised} support is
\[
\operatorname{supp}(z_d) = \{0\} \ \text{at } \kappa = 0,
\qquad
\{-\bar z,+\bar z\} \ \text{at } \kappa = 1,
\qquad
\{-\bar z,0,+\bar z\} \ \text{at every } \kappa \in (0,1).
\]
There is no conflict with hypothesis (P-2): the set $\{-\bar z,0,+\bar z\}$ is the noise
\emph{alphabet}, finite at every $\kappa$ and the only thing Step~3's finiteness argument needs,
while $\operatorname{supp}(z_d)$ is the subset of that alphabet carrying positive probability at the
maintained $\kappa$, and it is the object over which a zero-probability argument must be quantified.
Call a pooled history $\mathcal H_d^P$ \emph{reachable} if it carries strictly positive probability
under some plan and a positive-probability set of signals: there are $j \in \mathcal J$ and a Borel
set $S$ with $\Prb(s \in S) > 0$ such that $X_{d'}-q_{jd'}(s) \in \operatorname{supp}(z_{d'})$ for
every $d' \le d$ and every $s \in S$, with the flag coordinate agreeing with $\ind\{f_j(s) \le d\}$.
Reachability so defined is a property of the menu and the noise law alone: it does not depend on $k$,
and it does not depend on the perturbation stage. The whole-history form is what the argument needs,
since a history each of whose marks is individually attainable under \emph{some} plan need not be
attainable under any \emph{one} plan.

\emph{(b) The limit exists at every reachable history, and it is a limit of the joint posterior.}
Index the perturbation by $n$: at stage $n$ every signal type plays every plan $j \in \mathcal J$
with weight at least $1/n$, the remaining mass following $j_k$ of \eqref{eq:P1-jk}. Writing
\begin{equation}
\label{eq:P1-wn}
w_n(j \mid s) \;=\; (1-t_n)\,\ind\{j = j_k(s)\} + \frac{t_n}{J},
\qquad
t_n := \frac{J}{n} \downarrow 0,
\end{equation}
for the stage-$n$ mixing weight --- the perturbation mass is written $t_n$ rather than with a Greek
letter because $\varepsilon$ is the signal noise of Definition~\ref{def:primitives} --- and
\[
L_j\bigl(\mathcal H_d^P \mid s\bigr)
= \prod_{d' \le d} \Prb\bigl(z_{d'} = X_{d'}-q_{jd'}(s)\bigr)
\cdot \ind\bigl\{\text{flag coordinate} = \ind\{f_j(s) \le d\}\bigr\}
\]
for the likelihood of the history of (a), the stage-$n$ joint density over $(j,s)$ at that history is
\begin{equation}
\label{eq:P1-mun}
\mu_n(j,s) \;=\; \frac{w_n(j \mid s)\,L_j(\mathcal H_d^P \mid s)\,\varphi_s(s)}
{\sum_{j' \in \mathcal J} \int w_n(j' \mid s')\,L_{j'}(\mathcal H_d^P \mid s')\,\varphi_s(s')\,
\mathrm{d}s'}.
\end{equation}
By hypothesis (P-2) the plan menu is finite and by Step~3 the pooled history alphabet is finite, so
numerator and denominator are finite sums of terms polynomial in $1/n$ with coefficients not
depending on $n$. At a reachable history the denominator is strictly positive for every $n$: the
witnessing pair $(j,S)$ contributes at least $(1/n)\int_S L_j(\mathcal H_d^P \mid s)\varphi_s(s)\,
\mathrm{d}s$, which is strictly positive because each factor is positive on $S$ by the definition of
reachability and, the mark $q_{jd'}(\cdot)$ taking finitely many values, $L_j$ takes finitely many
positive values on $S$ and is bounded below by their minimum; every other term is nonnegative. A
ratio of polynomials in $1/n$ whose denominator is nonzero for all large $n$ converges as
$n \to \infty$, pointwise in $(j,s)$, and the limit is the plan-uniform-weighted joint law restricted
to the history. This is where the finiteness clauses pay: with a continuum of pooled histories the
limit would need a separate argument.

The joint form is what the step must deliver rather than a flourish. Step~4 shows that the pricing
map reads the information set through $(\hat v,\pi)$; while $\pi$ is a functional of the posterior
over plans alone, $\hat v(\mathcal I) = \Eop[v \mid \mathcal I]$ is a functional of the posterior over
\emph{signals}. Integrating the signal out first would deliver $\pi$ and leave $\hat v$ undefined.
Passing from $\mu_n$ to the belief, $\hat v_n = \sum_j \int \bigl(\mu_v+\beta(s-\mu_v)\bigr)
\mu_n(j,s)\,\mathrm{d}s$, and the envelope needs a case split on the denominator. Write
\begin{equation}
\label{eq:P1-Zn}
Z_n \;=\; (1-t_n)\,\Lambda_k + \frac{t_n}{J}\,\Lambda_u
\end{equation}
for the denominator of \eqref{eq:P1-mun} with $w_n$ expanded, where
\[
\Lambda_k = \int L_{j_k(s')}\bigl(\mathcal H_d^P \mid s'\bigr)\varphi_s(s')\,\mathrm{d}s',
\qquad
\Lambda_u = \sum_{j'} \int L_{j'}\bigl(\mathcal H_d^P \mid s'\bigr)\varphi_s(s')\,\mathrm{d}s'
\]
are the unperturbed and the plan-uniform likelihood aggregates. If $\Lambda_k > 0$ --- the history is
on path under
$k$ --- then $Z_n \ge \Lambda_k/2$ for all large $n$, and
$2\lvert \mu_v+\beta(s-\mu_v)\rvert\varphi_s L_j/\Lambda_k$ is an integrable envelope by the Gaussian
tail of clause (TR-i) and the integrability clause of Assumption~\ref{asm:A2p}, so dominated
convergence applies. If $\Lambda_k = 0$ --- the $k$-null case --- then the $(1-t_n)$ term vanishes
$\Phi_s$-almost everywhere in numerator and denominator alike, so $\mu_n = L_j\varphi_s/\Lambda_u$
exactly and free of $n$, and there is nothing to pass to the limit; note that this limit does not
depend on $k$, which is what makes one family serve at every $k \in \Theta$. Without the split the
bare claim would be false as stated, since $\mu_n \le \varphi_s L_j/Z_n$ with $Z_n \downarrow 0$ is
not a uniform envelope. Either way $\hat v_n \to \hat v_\infty$ and, by the $1$-Lipschitz bound
\eqref{eq:P1-dP}, the prices converge with the beliefs. The limiting belief and the limiting price
therefore exist at every reachable pooled history, on path and off, and on path the belief agrees
with the Bayes posterior. This is load-bearing where it is least obvious: Step~13 evaluates the
payoff to \emph{every} plan in the menu, including plans carrying zero probability under $k$ on a
collapse face, whose pooled-execution bracket reads prices at $k$-null histories --- so those prices
must be defined, not merely constrained.

\emph{(c) Unreachable histories carry no requirement, and a convention keeps the payoff defined at
every signal.} An unreachable history has probability zero under \emph{every} plan profile,
perturbed or not. It is null under nature rather than off the players' path, so requirement~(vi) of
Definition~\ref{def:cpbe} asks nothing of it and requirement~(iv) prices nothing there. Almost
everything downstream is already clear of such histories: Step~11's pooled-execution bracket
integrates $P_d^P(\mathcal H_d^P)$ against the law of $z_{0:H}$ under the plan actually played, which
for $\Phi_s$-almost every $s$ puts mass only on reachable histories, and the same holds for every
deviation in the round-2 action set of hypothesis (P-9), since those share the same pooled path.

The exceptional signals must nonetheless be named. Reachability in (a) requires a
positive-probability signal set, while the joint mark-and-flag vector
$\bigl((q_{jd'}(s))_{d' \le d},\,\ind\{f_j(s) \le d\}\bigr)$ takes finitely many values by hypothesis
(P-2), so its level sets partition the line into finitely many Borel cells --- and a particular cell
may be $\Phi_s$-null while being nonempty. At such a signal, plan $j$'s own realised pooled histories
are unreachable and Step~11's bracket would read a price that (b) has not defined. We therefore fix
once and for all a \emph{reference belief} $(\hat v_\circ,\pi_\circ)$, for definiteness the prior pair
$\bigl(\mu_v,\Prb(a=1)\bigr)$, and assign to every unreachable pooled history the inner root
$P_\circ$ at that belief, the unique solution of
$\mathcal P_{(\hat v_\circ,\pi_\circ)}(P_\circ) = P_\circ$, which exists and is unique by Step~7 under
hypothesis (P-11). The payoff is then defined at every triple $(j,s,k)$, and no requirement of
Definition~\ref{def:cpbe} is touched, since requirements (iii), (iv) and (vi) constrain beliefs and
prices only at histories carrying positive probability under some profile.

The form of the convention matters, and the unconditional mean $\Eop[Y]$ will not serve. That
quantity is an unconditional average of realised control-node values and is in general not a root of
$\mathcal P_{\mathcal I}$ at any belief, so the constructed object would carry prices that are not
inner fixed points at nodes where the blockholder does trade, while the proposition's conclusion says
``prices at their inner fixed points'' without qualification. The precedent of
$P_{-1}^P = \Eop[Y]$ in Definition~\ref{def:prices} does not transfer: that is the pre-trading node, a
different object. With the reference-belief root, every price in the constructed object is an inner
fixed point at the belief carried there. The choice is not innocuous in one narrow respect, and we do
not pretend otherwise: it can change the payoff on a $\Phi_s$-null set of signals, and hence can move
a component of the outer map, which is an infimum over a pointwise condition. What follows is that
the proposition is an existence statement about the object built from this fixed convention; a
different admissible convention yields another equilibrium, not a failure of this one.

\emph{(d) The boundary values of $\kappa$.} At $\kappa \in \{0,1\}$ the realised support degenerates
as displayed in (a): to $\{0\}$ and to $\{\pm\bar z\}$ respectively. A pooled history needing a mark
outside the surviving support is null under \emph{every} plan profile and at every perturbation
stage, so it is off nature's path rather than off the players', it carries no requirement under
Definition~\ref{def:cpbe}\,(vi), and no step above or below reads it. Parts (a) to (c) are quantified
over $\operatorname{supp}(z_d)$ and so cover both endpoints in one sentence; the claim therefore
holds on the full domain $\kappa \in [0,1]$, with no cut to $\kappa \in [0,1)$ anywhere. This is the
extension route of clause~(ii) of the proposition rather than a restriction. What the endpoints
change is \emph{which} histories are reachable, never whether the reachable ones carry a limit; and
the correction runs in the direction of honesty rather than retreat, since the earlier claim that a
limiting belief exists at \emph{every} pooled history is false at $\kappa = 1$ and is withdrawn.
Nothing here asserts that the flagged cell carries positive probability at $\kappa = 1$, and nothing
below does either.

\smallskip
\noindent\textbf{Step 10 (the flagged belief is the point mass at the generating pair, and it is a
version).}
By hypothesis (P-6), Assumption~\ref{asm:A7J} (A7-J), the map $(j,s) \mapsto \sigma_F$ is injective on the
flagged-pair set, so each flagged tuple in its image is generated by exactly one pair, and the Borel
inverse $\iota_F$ of Step~6(c) returns it.

The version must be stated explicitly, because the naive argument is unavailable. The signal is
Gaussian by clause (TR-i), so the pair $(j,s)$ carries probability zero under every stage-$n$
perturbation --- what carries positive weight is the plan conditional on the type. A conditional law
given $\sigma_F$ is therefore defined only up to $\Phi_s$-null sets, and the sentence ``that pair has
strictly positive weight, so the perturbed posterior at $\sigma_F$ places probability one on
$\iota_F(\sigma_F)$'' would apply a positive-probability argument to a null event. What is true, and
is all that is needed, is this: because $\iota_F$ is a genuine pointwise Borel map,
$\delta_{\iota_F(\sigma_F)}$ is a \emph{version} of the regular conditional law given $\sigma_F$ at
every stage $n$ --- the defining disintegration identity holds tuple by tuple, since the conditioning
$\sigma$-field separates the generating pairs --- and it is invariant in $n$, hence its own limit.
This proof selects that version at every tuple in the image; any $\Phi_s$-almost-everywhere-equal
version satisfies requirements (iii) and (vi) of Definition~\ref{def:cpbe} equally well, and nothing
below distinguishes them. So the belief is available and forced up to null sets, not unique.
Consequently the flagged belief is $\hat v(\sigma_F) = \mu_v + \beta\bigl(\iota_F(\sigma_F)_s -
\mu_v\bigr)$ at every image tuple, on path and off, and Step~6's family is simultaneously the on-path
Bayes family and the off-path limit family.

What ``off path'' covers here must be said precisely. It covers flagged tuples generated by pairs
$(j,s)$ that the conjecture $k$ does not select --- these are exactly the pairs
Assumption~\ref{asm:A7J} (A7-J) reaches and its on-path form, Assumption~\ref{asm:A7p} (A7$'$), does not. It does
\emph{not} by itself cover tuples outside the \emph{image} of $(j,s) \mapsto \sigma_F$, the tuples a
round-2 deviation to an off-menu order would produce, and no step assigns those a belief. The step
stands here on hypothesis (P-9), the definitional round-2 action-set hypothesis: because the round-2
action set \emph{is} the plan-generated set, no such tuple arises, and the image exhausts the flagged
tuples that can be reached. Off-path beliefs at flagged nodes therefore carry no free parameter,
which is a consequence of Assumption~\ref{asm:A7J} (A7-J) worth recording, since it removes the usual
arbitrariness at exactly the nodes the disclosure mechanism runs through. By hypothesis (P-1) the
pair $(v,\xi)$ remains conditionally independent of the pooled residual given the signal, so nothing
further is needed to price the node.

\medskip
\noindent\emph{Part D: sequential optimality.}

\smallskip
\noindent\textbf{Step 11 (there are exactly two decision points, and the flagged continuation is
deterministic given $(j,s)$).}
Section~\ref{sec:timing} places the plan choice at date $0$ and, when $D = 1$, the flagged order in
round~2, with no re-optimisation in between by hypothesis (P-8). There is thus no pooled decision
node after date $0$: requirement~(ii) of Definition~\ref{def:cpbe}, read on the pooled component, is
satisfied by the timing itself rather than by an argument, and the only genuine sequential-optimality
requirement is the round-2 order. Two features of the decomposition below come from the
flag-termination clause of Assumption~\ref{asm:flagterm}, the second half of hypothesis (P-8), and
not from the no-feedback clause: that the pooled execution runs over $d \le f_j$ and stops there, and
that $Q_j^F = b_j^*(s)-B_j^F(s)$ is the blockholder's \emph{whole} residual position.

On the flagged branch, by Step~2 the objects $B_j^F(s)$ and $Q_j^F(s)$ are deterministic in $(j,s)$;
by Step~6 the flagged price is a function of $\sigma_F$ alone and hence deterministic in $(j,s)$; and
by Definition~\ref{def:prices} the control-node price on that branch is $P^F$. Write $G_j(s;k)$ for
the flagged \emph{trading terms} and $E_j(s;k)$ for the pooled-execution expectation,
\begin{align*}
G_j(s;k) &\;=\; b_j^*(s)\,\Eop\bigl[Y \mid s,j,D=1\bigr] - P^F(\sigma_F)\,Q_j^F(s),\\
E_j(s;k) &\;=\; \Eop_{z}\Bigl[\textstyle\sum_{d \le f_j} P_d^P\bigl(\mathcal H_d^P\bigr)
\bigl(B_j(s,d)-B_j(s,d-1)\bigr)\Bigr],
\end{align*}
the second expectation being over the noise alone. The objective of hypothesis (P-12),
Definition~\ref{def:U} with the display \eqref{eq:Uj}, then splits as
\begin{equation}
\label{eq:P1-Udecomp}
U_j(s;k) \;=\; \underbrace{G_j(s;k) - a_j C_j(s)}_{\text{flagged continuation: deterministic in }(j,s)}
\;-\; \underbrace{E_j(s;k)}_{\text{pooled execution: fixed by the pooled path}} .
\end{equation}
On the flagged branch $a_j = 1$ by
Assumption~\ref{asm:A4}, so the engagement term is $C_j(s)$; the flag factor is carried to match
\eqref{eq:Uj} and changes nothing in Steps~12 and~13, since those run on the flagged set. By
Step~9(c), for $\Phi_s$-\emph{almost every} signal every pooled history the second bracket weighs is
reachable, so the prices it reads are the ones Steps~5 and~9 supply; at the exceptional signals it
reads the reference-belief root of Step~9(c) instead, which is itself an inner root, and which is
what keeps the payoff defined at every signal for Step~13's pointwise argmax. The noise enters the
first bracket nowhere: $\Eop[Y \mid s,j,D=1]$ depends on $(v,\xi)$ and on $P^F$, and $\xi$ is
independent of the noise by hypothesis (P-1) while $P^F$ is noise-free by Step~6.

\smallskip
\noindent\textbf{Step 12 (the flagged component is sequentially optimal at \emph{every} flagged pair,
selected or not: price invariance, cancellation, and the continuation-cost clause).}
This is the step that discharges requirement~(ii) of Definition~\ref{def:cpbe} on the flagged
component, and it does so without appealing to date-$0$ optimality --- which is precisely what lets
it cover flagged pairs that the cutoff vector does not select. Such pairs are genuine round-2
information sets: a date-$0$ deviation to a non-selected plan that flags creates one, and
requirement~(ii) binds there exactly as it binds on the selected plan.

Fix a flagged pair $(j,s)$, with no assumption that $j = j_k(s)$, and let $j'$ range over the class
generating the round-2 action set of hypothesis (P-9),
\[
\mathcal Q_j(s) \;:=\; \bigl\{Q_{j'}^F(s) : j' \in \mathcal J \text{ shares } j\text{'s pooled path up
to } f_j(s) \text{ and } a_{j'} = a_j\bigr\},
\]
which is what the flagged-round action set \emph{is} under that hypothesis. The shared pooled path
forces $c_{j'}(s) = c_j(s)$, since the crossing date is a first-hitting index of a path the two
plans share and $c_j \le f_j - T \le f_j$; hence $f_{j'}(s) = f_j(s)$ and $B_{j'}^F(s) = B_j^F(s)$ by
Step~2 and hypothesis (P-7), and $a_{j'} = a_j = 1$ on the flagged set by Assumption~\ref{asm:A4}.
The deviation's flagged tuple is therefore $\sigma_F(j',s) = \bigl(B_j^F(s),Q_{j'}^F(s),1\bigr)$ ---
in the image of the flagged-pair map, where Step~10 applies --- and it differs from
$\sigma_F(j,s)$ in the $Q^F$ coordinate alone. Two remarks on the class before the argument. The
sharing relation is an equivalence relation on the flagged set at each signal, since the mutual
agreement just derived makes it symmetric; and it is a genuine pair only on menus carrying two or
more Voice plans that share a pooled path, being the singleton $\{j\}$ on any single-Voice menu,
where Exit and Hold never cross $\tau$ under $b_0 < \tau$.

\emph{(a) The flagged price does not move across the class.} By Step~10 the belief at
$\sigma_F(j',s)$ is the point mass at $(j',s)$, so
$\hat v\bigl(\sigma_F(j',s)\bigr) = \mu_v + \beta(s-\mu_v)$ --- a function of the signal alone, the
same for every member of the class --- and $\pi = 1$ on the flagged cell by Step~6(a), so
$\bar m = m_1$. By Step~4 the inner pricing map depends on the information set only through
$(\hat v,\pi)$. Hence
\begin{equation}
\label{eq:P1-PFinv}
P^F\bigl(\sigma_F(j',s)\bigr) \;=\; \mathcal G_F\bigl(\hat v(s)\bigr) \;=:\; P^F(s)
\qquad \text{for every } j' \text{ in the class,}
\end{equation}
so the round-2 order carries no price impact across the menu. It is Assumption~\ref{asm:A7J} (A7-J) that
pins the belief at the same signal for every class member, and it is uniqueness of the inner root
from Step~7 that makes $P^F(s)$ a number rather than a selection. The argument at this display is
where the joint form of the injectivity hypothesis is indispensable and its on-path form,
Assumption~\ref{asm:A7p} (A7$'$), is not enough: the tuples in question are generated by pairs the cutoff
vector need not select.

\emph{(b) At a flagged node the blockholder values a share at exactly $P^F(s)$.} The blockholder
knows $(j,s)$ and the realised pooled history; given $(j,s)$ the latter is a function of
$z_{0:f_j}$ by Step~2, and the noise is independent of $(v,\varepsilon,\xi)$ by hypothesis (P-1), so
it carries no information about the terminal payoff:
$\Eop\bigl[Y \mid s,j',\mathcal H_{f^-}^P,D=1\bigr] = \Eop\bigl[Y \mid s,j',D=1\bigr]$. By the same
hypothesis $\mathsf B$ is a function of $\xi$ alone given the $\sigma_F$-measurable $P^F$, so with
$p = p\bigl(P^F(s)\bigr)$ from Step~4,
\[
\Eop\bigl[Y \mid s,j',D=1\bigr]
\;=\; (1-p)\bigl(\Eop[v \mid s] + \Delta_V\bigr) + p\bigl(P^F(s)+m_1\bigr).
\]
Assumption~\ref{asm:A7J} (A7-J) makes the market's flagged posterior the blockholder's own: by (a) we
have $\hat v\bigl(\sigma_F(j',s)\bigr) = \Eop[v \mid s]$, so the right-hand side above is
$\mathcal P_{\mathcal I}\bigl(P^F(s)\bigr)$ evaluated at the flagged information set, which equals
$P^F(s)$ because $P^F(s)$ solves the inner fixed point by Steps~4 and~6. Hence
\begin{equation}
\label{eq:P1-EYflag}
\Eop\bigl[Y \mid s,j',\mathcal H_{f^-}^P,D=1\bigr] \;=\; P^F(s).
\end{equation}
No informational rent survives into round~2: full separation is what Assumption~\ref{asm:A7J} (A7-J) buys,
and \eqref{eq:P1-EYflag} is what it costs.

\emph{(c) The flagged-order terms cancel.} The flag terminates the pooled round by
Assumption~\ref{asm:flagterm}, so at round~2 the pooled execution is complete, sunk, and common to
every class member --- it is the bracket $E_j(s;k)$ of \eqref{eq:P1-Udecomp}, which the shared path
makes identical across the class for every noise draw. Write $V(j')$ for the round-2 continuation
value of submitting $Q_{j'}^F(s)$: the terminal position valued at the control node, less what the
flagged order costs, less the engagement cost the deviator bears. Then
\begin{equation}
\label{eq:P1-Vcancel}
\begin{aligned}
V(j')
&\;=\; \bigl(B_j^F(s)+Q_{j'}^F(s)\bigr)\,\Eop[Y \mid \cdot\,]
\;-\; P^F(s)\,Q_{j'}^F(s) \;-\; (\text{engagement cost})\\[2pt]
&\;=\; B_j^F(s)\,P^F(s) \;-\; (\text{engagement cost}),
\end{aligned}
\end{equation}
using \eqref{eq:P1-EYflag} for the valuation and \eqref{eq:P1-PFinv} for the price at the deviation
tuple. Every appearance of $Q_{j'}^F$ has cancelled, and with it every appearance of $b_{j'}^*$: the
class member's identity survives only through the engagement cost. In the notation of
\eqref{eq:P1-Udecomp} the same computation reads $G_{j'}(s;k) = B_j^F(s)P^F(s)$ for every class
member, so on flagged plans $U_{j'}(s;k) = B_j^F(s)P^F(s) - C_{j'}(s) - E_j(s;k)$.

\emph{(d) The continuation-cost clause closes the step, and the conclusion is convention-free.}
Display \eqref{eq:Uj} does not date the engagement cost, and Remark~\ref{rem:Ctiming} records the two
available readings. Under the \emph{plan-completion} reading, submitting $Q_{j'}^F(s)$ is completing
plan $j'$, so the deviator bears $C_{j'}(s)$; under the \emph{sunk} reading, the filing has landed
and $D = 1 \Rightarrow a = 1$ is public by Assumption~\ref{asm:A4}, so the engagement cannot be
unmade and the deviator bears $C_j(s)$ whatever order is submitted. Under the sunk reading
\eqref{eq:P1-Vcancel} is constant on the class outright. Under the plan-completion reading,
hypothesis (P-10) --- continuation-cost equivalence on this same set, $C_{j'}(s) = C_j(s)$ for every
class member --- makes it constant again. Either way
\[
V(j') \;=\; B_j^F(s)\,P^F(s) - C_j(s)
\qquad \text{for every } j' \text{ in the class,}
\]
so every element of $\mathcal Q_j(s)$, the specified $Q_j^F(s)$ included, attains the maximum. The
flagged component is sequentially optimal at every flagged pair $(j,s)$, on the selected plan and off
it, which is requirement~(ii) of Definition~\ref{def:cpbe} in full on its flagged half.

Where hypothesis (P-10) bites is worth stating exactly, since it is vacuous over a large part of its
range. Under the sunk reading it is not consumed at all. Under the plan-completion reading, at a
\emph{selected} $j$ date-$0$ optimality would already do the work: by (c) the within-class comparison
$U_j \ge U_{j'}$ \emph{is} $C_j \le C_{j'}$, which defeats the deviation. But at a flagged pair the
cutoff vector does not select there is no date-$0$ optimality to appeal to, and without the clause the
deviator strictly prefers the class member with the smallest engagement cost. So the clause's bite is
precisely the flagged nodes off the equilibrium plan under the plan-completion reading, and what it
buys is that the conclusion holds under both readings without this proof adjudicating a silence in
the model's own accounting. It is stated as an equality rather than one-sidedly for two reasons: the
sharing relation is symmetric, as shown above, so a clause imposed at every pair of a class in one
direction is imposed in both; and the clause cannot be indexed by the equilibrium's selection, since
Brouwer's theorem does not say which fixed point it returns and the requirement lands at pairs no
cutoff vector selects, where ``the selected plan'' names nothing.

Two things this step does \emph{not} establish should be recorded here rather than left to be
inferred. First, a class on which the \emph{trading} terms differ cannot be built. A witness fixing
$G_{j'}(s;k) - G_j(s;k) = \delta > 0$ across a class is inconsistent with requirements (iv) and (vi)
of Definition~\ref{def:cpbe} in force: by (a) and (b), at pinned beliefs and inner-fixed-point prices
$G_{j'} = B_j^F(s)P^F(s)$ for \emph{every} class member, so $\delta = 0$ necessarily. Such a witness
fixes as a primitive what equilibrium determines, and is available only where the flagged price is
not the fixed point of \eqref{eq:P1-inner} or the flagged belief is not the one
Assumption~\ref{asm:A7J} (A7-J) pins. The live wedge is in the engagement cost, which is where equilibrium
leaves room for one, and that is what hypothesis (P-10) pays for. Second, the converse direction is
about hypothesis (P-9) and not about the argument above. Without it --- if round~2 offered the full
interval $[0,\bar b-B_j^F(s)]$ --- an order outside the plan-generated set would produce a flagged
tuple outside the image of $(j,s) \mapsto \sigma_F$, where Step~10 pins nothing and no step assigns a
belief; the deviation's price, and with it the whole comparison, would be undefined until a belief
were supplied, and the supplied belief would decide the answer. By (c) the on-image flagged payoff is
$B_j^F(s)P^F(s)-C_j(s)$ with the order cancelled out, so an off-image deviation is profitable or not
entirely according to how the assigned off-image belief compares with $\Eop[v \mid s]$, and Step~9's
plan-only perturbation constrains that choice at no stage $n$. Sequential optimality of the flagged
component is therefore not a free consequence of complete contingent plans, and it does not follow
from Assumptions~\ref{asm:A1} through \ref{asm:A7J} (A1 through A7-J) on their own: hypothesis (P-9) is \emph{a} sufficient
condition that delivers it, and it is a restriction on the round-2 action set rather than on the
menu. Strengthening the injectivity hypothesis does not change this, since the menu witnessing the
failure may be taken with $b^*$ strictly increasing on the whole line, so that
Assumption~\ref{asm:A7J} (A7-J) holds on it and requirement~(ii) still fails. Whether hypothesis (P-9) is
the \emph{weakest} such condition is not claimed and is open: the textbook route to sequential
rationality at an unreached node is not a restriction on the action set but a choice of off-path
beliefs, and Definition~\ref{def:cpbe}\,(vi) leaves round-2 orders outside the menu image reached at
no stage $n$, so their limit beliefs are unconstrained by the perturbation used here. Whether some
admissible choice deters every off-menu deviation is a genuine question --- by (c) an off-image
deviation to $Q'$ at belief $\hat v'$ earns
$B_j^F\,\Eop[Y \mid \cdot\,] + Q'\bigl(\Eop[Y \mid \cdot\,]-P'\bigr)$ less the cost, so the assigned
belief moves the deviation's gain and the incumbent's own valuation at once, in opposite directions
--- and no step of this proof addresses it.

\medskip
\noindent\emph{Part E: the outer map and Brouwer.}

\smallskip
\noindent\textbf{Step 13 (a well-defined weakly ordered best-response map; the ordering and self-map
content of Assumption~\ref{asm:A6} (A6) is a consequence, not an assumption).}
Fix $k \in \Theta$. Steps~5, 6, 9 and~10 determine the pooled and flagged price families and the
belief system, and Step~11 then determines $U_j(s;k)$ for every plan and \emph{every} signal, by the
convention of Step~9(c); each value is finite by hypothesis (P-2), which supplies local boundedness
in $(s,\vartheta)$ pointwise and integrability in expectation. By the second clause of
Assumption~\ref{asm:A3}, hypothesis (P-3), the preferred plan is weakly increasing in the signal, so
the set $\mathcal S(k)$ of \emph{weakly increasing selections} from
$s \mapsto \arg\max_{j \in \mathcal J} U_j(s;k)$ is nonempty. The selection is named rather than
merely asserted to exist:
\begin{equation}
\label{eq:P1-jstar}
j^\star(\cdot\,;k) \;:=\; \text{the largest element of } \mathcal S(k),
\qquad
j^\star(s;k) = \max\bigl\{\mathfrak w(s) : \mathfrak w \in \mathcal S(k)\bigr\}.
\end{equation}
It exists and is itself a weakly increasing selection. \emph{Closure under pointwise maximum:} for
$\mathfrak w_1,\mathfrak w_2 \in \mathcal S(k)$ the map $s \mapsto \max\{\mathfrak w_1(s),
\mathfrak w_2(s)\}$ takes at each signal one of the two values, both in the argmax there, so it is a
selection, and a pointwise maximum of two weakly increasing maps is weakly increasing.
\emph{Attainment:} the menu is finite by hypothesis (P-2), so the pointwise supremum is a maximum;
fix $s$ and let $\mathfrak w \in \mathcal S(k)$ attain it there. Then $j^\star$ is a selection, since
$j^\star(s;k) = \mathfrak w(s)$ lies in the argmax at $s$, and it is weakly increasing, since for
$s < s'$ we have $j^\star(s';k) \ge \mathfrak w(s') \ge \mathfrak w(s) = j^\star(s;k)$. So
$j^\star(\cdot\,;k)$ is canonical, single-valued and monotone.

The largest \emph{maximiser} would not do, and the counterexample sits inside
Assumption~\ref{asm:A3}: let $U_2-U_1 \le 0$ everywhere with equality at exactly one signal $s_0$, a
tangential touch, so that there are no crossings and the constant selection $j \equiv 1$ is weakly
increasing, whence both clauses of that assumption hold. The largest maximiser is then $1,\dots,1,
2,1,\dots,1$: it is not weakly increasing, the set $\{s : j^\star \ge 2\} = \{s_0\}$ is not an upper
set, and no cutoff vector represents it, which would break the assembly in Step~17(i). The largest
weakly increasing selection is $j \equiv 1$ there, as it should be. Define now
\begin{equation}
\label{eq:P1-Tmap}
\Tmap_i(k;\vartheta) \;=\; \inf\bigl\{s \in [\underline s,\overline s] :
j^\star(s;k) \ge i+1\bigr\},
\qquad i = 1,\dots,J-1,
\qquad \inf\emptyset := \overline s .
\end{equation}
Since $\{s : j^\star \ge i+2\} \subseteq \{s : j^\star \ge i+1\}$, the infima satisfy
$\Tmap_1(k) \le \Tmap_2(k) \le \dots \le \Tmap_{J-1}(k)$, and every component lies in
$[\underline s,\overline s]$ by construction. The corner convention is the display's
$\inf\emptyset := \overline s$: a plan optimal nowhere contributes an empty upper set, its cutoff
sits at the top of the bracket, and it simply never appears in the range of $j^\star$. Because
$j^\star(\cdot\,;k)$ is weakly increasing, each set $\{s : j^\star \ge i+1\}$ \emph{is} an upper set,
so these infima genuinely represent it: $j^\star(\cdot\,;k)$ agrees with the induced map
$j_{\Tmap(k)}$ of \eqref{eq:P1-jk} at every signal except possibly the finitely many $\Tmap_i(k)$ at
which the upper set fails to contain its own infimum, where the two differ by the boundary convention
alone, and requirement~(i) of Definition~\ref{def:cpbe} pins no convention there. Hence
$\Tmap(k;\vartheta) \in \Theta$ for every $k \in \Theta$, including on the collapse faces where
consecutive components coincide and the corresponding plan carries zero probability --- which the
weak inequalities of Definition~\ref{def:cpbe} admit. The ``weakly ordered'' and ``maps $\Theta$ into
itself'' halves of Assumption~\ref{asm:A6} are thus derived here from the monotone-preferred-plan
clause of Assumption~\ref{asm:A3}. What remains genuinely assumed in hypothesis (P-5) is the bracket
$[\underline s,\overline s]$ and the continuity of $\Tmap$; Steps~14 and~15 take those in turn. The
tie-break and the corner convention named in \eqref{eq:P1-jstar} and \eqref{eq:P1-Tmap} are the named
selection the proposition's second trailing paragraph requires before
Assumption~\ref{asm:A6} can be applied at all, since a correspondence cannot be called continuous.

\smallskip
\noindent\textbf{Step 14 (the bracket: derivable in the four-action specialisation, assumed at the
level of generality the model is stated at).}
In the four-action version of this model, the bracket is proved rather than assumed. There the
blockholder's payoff to each action is affine in the posterior mean with intercepts bounded uniformly
over conjectures --- prices lie in a bounded interval and the entry probability lies in $[0,1]$ ---
and with totally ordered slopes: zero for Exit, one for Hold and Quiet Voice, and strictly more than
one for Public Voice, while the engagement cost is continuous, strictly positive and strictly
decreasing with full range on the half-line. Each adjacent indifference condition then equates two
affine functions whose slopes differ, except for the Hold and Quiet Voice pair, where the slopes tie
and the comparison reduces to a strictly decreasing cost schedule meeting a bounded constant. Every
indifference signal is therefore finite and bounded uniformly in the conjecture, and the union over
the finitely many adjacent pairs gives one bracket serving all of them.

That argument uses the affine payoff form with ordered slopes, and nothing in this model's
hypotheses imposes it on a general finite menu. Assumption~\ref{asm:A3} imposes single crossing and a
monotone preferred plan only; Definition~\ref{def:U} fixes the accounting of the payoff without
restricting its shape in the posterior mean, and names only plan-locality and integrability as the
properties ever used, neither of which is a shape restriction. At this level of generality the common
bracket is therefore \emph{assumed}, and it is the first of the two things hypothesis (P-5) is doing.

\smallskip
\noindent\textbf{Step 15 (continuity: where hypothesis (P-5) assumes rather than derives, and what it
assumes).}
Continuity of $U_j(s;k)$ in $k$ at fixed $(j,s)$ is available. The pooled and flagged inner prices
enter the payoff, by Step~4, only through $(\hat v,\pi)$; those are ratios of integrals over signal
intervals with endpoints $k$ and are continuous in $k$ wherever the conditioning event has
probability bounded away from zero, and at histories of vanishing probability the perturbation limit
of Step~9 supplies the value, while the root map is $1$-Lipschitz in the belief by
\eqref{eq:P1-dP}. That is continuity in $k$ at fixed $(j,s)$, \emph{and it is not enough}: continuity
in $k$ at fixed $s$ together with continuity in $s$ at fixed $k$ is strictly weaker than continuity
in the pair, and it is the joint statement that the crossing-point argument below consumes.
Continuity of $\Tmap$ needs two further things that no hypothesis of
Section~\ref{sec:hypotheses} supplies.

\begin{enumerate}[label=(\roman*),leftmargin=2.4em,itemsep=3pt]
\item \emph{Joint continuity.} The map $(s,k) \mapsto U_j(s;k)$ is continuous on
  $[\underline s,\overline s] \times \Theta$ for each plan. This is stated as a condition, not
  inferred from the two separate continuities. It is plausible from the structure --- finitely many
  plans, the inner root $1$-Lipschitz in the belief, and $(\hat v,\pi)$ ratios of integrals over
  signal intervals with endpoints $k$ --- and what it needs in the signal direction is continuity of
  $s \mapsto \bigl(B_j(s,\cdot),b_j^*(s),C_j(s)\bigr)$. Clause (TR-ii) imposes monotonicity on the
  stake path and Borel regularity, and nothing more; a plan acquiring a block discontinuously at a
  signal trigger is permitted and makes the payoff jump, at which point the best-response cutoff is a
  jump point rather than a crossing point and moves discontinuously with $k$.
\item \emph{Transversality.} For every adjacent pair $(i,i+1)$ and every $k \in \Theta$, the
  indifference set $\{s : U_{i+1}(s;k) = U_i(s;k)\}$ has empty interior. Assumption~\ref{asm:A3} says
  the difference crosses zero at most once, which does not exclude an interval on which it vanishes
  identically; on such an interval the cutoff is indeterminate, and as the interval opens and closes
  with $k$ the selection jumps.
\end{enumerate}

Under (i) and (ii), continuity of $\Tmap$ follows from the joint continuity together with the strict
sign change of $U_{i+1}-U_i$ at each crossing: the sign change locates the crossing and the joint
continuity moves it continuously with $k$. That is a topological argument rather than a calculus one;
the implicit-function theorem is the wrong tool here, since it would need the payoff differentiable
in $(s,k)$ and no hypothesis supplies that. Hypothesis (P-5) assumes the conclusion instead, by
asserting continuity of $\Tmap$ directly. That is the single largest assuming-rather-than-deriving
step in this proof, and (i) together with (ii) is the weakest pair of conditions we can name that
would replace it. Note also that (i) is not independent of hypothesis (P-6): a stake path flat on a
signal interval destroys injectivity there, so continuity cannot be bought by weakening monotonicity.

Two configurations in which the named sufficient conditions fail are worth recording, because they
show hypothesis (P-5) doing more work than Step~15 makes it look like it is doing. An engagement cost
constant on a signal interval, with the conjecture such that $U_{i+1}(\cdot;k)-U_i(\cdot;k)$ vanishes
there, makes every cutoff in that interval a best response, and the selection jumps as the plateau
opens and closes with $k$: condition (ii) fails. Worse, plateaus are structural on exactly the menus
hypothesis (P-10) exists for. By (c) of Step~12,
$U_{j'}(s;k) = B_j^F(s)P^F(s)-C_j(s)-E_j(s;k)$ is the \emph{same function of the signal} for every
member of a deviation class --- the shared pooled path gives a common stake at filing and a common
execution bracket, and hypothesis (P-10) gives a common cost --- so on any menu carrying two
\emph{adjacent} class members the difference $U_{i+1}-U_i$ vanishes identically on the whole flagged
region, and condition (ii) fails identically rather than on a knife-edge. Neither observation breaks
the proof: hypothesis (P-5) asserts continuity of $\Tmap$ outright, and the selection
\eqref{eq:P1-jstar} is single-valued across a plateau. But it should be said plainly that on
multi-Voice shared-path menus that hypothesis is being assumed at a configuration where its own named
sufficient condition provably fails. Remark~\ref{rem:A6record} carries the separate, measured record
that the continuity clause fails at the implemented calibration for the $\Theta$ this argument
constructs.

\smallskip
\noindent\textbf{Step 16 (Brouwer).}
By Step~1 and Definition~\ref{def:theta}, $\Theta$ is nonempty, compact and convex. By Step~13,
$\Tmap(\cdot\,;\vartheta)$ is a single-valued self-map of $\Theta$ under the largest weakly
increasing selection \eqref{eq:P1-jstar} and the corner convention of \eqref{eq:P1-Tmap}, and this is
the reading of Assumption~\ref{asm:A6} the step uses: that hypothesis asserts that $\Tmap$, so
selected, is a continuous self-map of $\Theta$. Brouwer's fixed-point theorem then supplies
$k^\star \in \Theta$ with $k^\star = \Tmap(k^\star;\vartheta)$. The fixed point may lie on a collapse
face, in which case the corresponding plan carries zero probability; the weak inequalities of
Definition~\ref{def:cpbe}\,(i) admit this.

\smallskip
\noindent\textbf{Step 17 (assembling the six requirements of Definition~\ref{def:cpbe}).}
Take $k^\star$ from Step~16 and check the definition item by item.

\begin{enumerate}[label=(\roman*),leftmargin=2.4em,itemsep=3pt]
\item \emph{Weakly ordered cutoff vector.} $k^\star \in \Theta$ by Step~16, and the equilibrium plan
  map is $j^\star(\cdot\,;k^\star)$ of \eqref{eq:P1-jstar} --- the largest weakly increasing selection
  from the pointwise argmax, weakly increasing by construction rather than by an appeal to
  Assumption~\ref{asm:A3} after the fact, and represented by $k^\star$ because a weakly increasing
  menu-valued map has upper sets for its upper level sets. It and the induced map $j_{k^\star}$ agree
  off the cutoff points and can differ \emph{at} them, since each $\Tmap_i$ is an infimum that need
  not be attained; requirement~(i) asks for a weakly ordered vector mapping the signal into a plan and
  does not pin the tie convention at the cutoffs, so taking the map to be $j^\star$ --- optimal at
  \emph{every} signal by construction --- is admissible and is what (ii) needs. Consistency with the
  conjecture the prices are built on is immediate: Steps~5, 6 and~9 price against $k^\star$, whose
  induced map is $j_{k^\star}$, and the two maps agree off the finitely many cutoff points, hence
  $\Phi_s$-almost surely, so every conditional probability, posterior and price is the same under
  either. The disagreement is invisible to (iii), (iv) and (v), which are statements about
  probabilities.
\item \emph{Sequentially optimal pooled and flagged components.} Pooled: there is no decision node
  after date $0$, by Step~11. Flagged: Step~12, under hypotheses (P-9) and (P-10), at every flagged
  pair, selected or not. Date-$0$ plan optimality: $k^\star$ is a fixed point of $\Tmap$ and the plan
  map is $j^\star(\cdot\,;k^\star)$, so $j^\star(s;k^\star) \in \arg\max_j U_j(s;k^\star)$ at every
  signal by the definition of the selection. No appeal to indifference at the cutoff points is
  needed, and none is available: indifference there would require continuity of
  $s \mapsto U_j(s;k)$, which is condition (i) of Step~15, explicitly not a hypothesis and explicitly
  not derived.
\item \emph{Bayes-consistent on-path beliefs.} Step~9 for pooled histories of positive probability
  under $k^\star$; Step~10 for flagged tuples, where Assumption~\ref{asm:A7J} (A7-J) supplies the point mass
  at $\iota_F(\sigma_F)$ as the selected version.
\item \emph{Competitive pooled and flagged prices at their fixed points.} Step~5 for the finite
  pooled family and Step~6 for the measurable flagged family, both solving
  $P(\mathcal I) = \Eop[Y \mid \mathcal I]$ by Step~4 --- at the beliefs of (iii) where the history
  carries positive probability under $k^\star$, and at Step~9(b)'s limit belief at the reachable
  histories that do not, since those are the histories the deviation payoffs defining $\Tmap$ read.
  At unreachable histories the requirement asks nothing and Step~9(c)'s convention supplies a value
  that is itself an inner root.
\item \emph{Bidder-entry rule.} The entry probability \eqref{eq:entry} is the one implied by the same
  pair $(P,\pi)$ at each control-node information set, by the derivation in Step~4.
\item \emph{Off-path beliefs as limits of full-support perturbations.} Steps~9 and~10: one family,
  fixed once in \eqref{eq:P1-wn} and used to define the price system at every $k \in \Theta$, at
  every pooled history reachable with positive probability under some plan profile; and, at flagged
  nodes, the version of Step~10 at every tuple in the image of the flagged-pair map, with no tuple
  outside that image arising under hypothesis (P-9).
\end{enumerate}

All six requirements hold, so the assembled object is a cutoff perfect Bayesian equilibrium, and the
construction has been carried out at an arbitrary $\kappa \in [0,1]$, endpoints included, by
Step~9(d). This proves the existence half of the proposition together with its clauses (i) to (v).

\smallskip
\noindent\textbf{Step 18 (a strengthening recorded here, not part of the proposition).}
Define instead the best-response correspondence
\[
\mathfrak T(k) \;=\; \bigl\{k' \in \Theta : k' \text{ represents some optimal weakly increasing plan
selection at } k\bigr\}.
\]
It is nonempty by Assumption~\ref{asm:A3}; its values are convex, because at an
indifference plateau the admissible values of a component form an interval and the ordering
constraints cut the product of those intervals by half-spaces; its values are compact, being closed
subsets of the compact $\Theta$; and its graph is closed by the maximum theorem, given that the
payoff is jointly continuous in $(s,k)$, which is exactly what condition (i) of Step~15 states as a
condition rather than deriving. Kakutani's fixed-point theorem then delivers a fixed point without
condition (ii) of Step~15 and without the continuity clause of Assumption~\ref{asm:A6}. It removes
the transversality condition and \emph{neither} condition (i) nor Step~14's bracket.
Definition~\ref{def:cpbe} fixes the Brouwer route for this proposition, so this is a remark and not a
part of the claim. \emph{Card note, 2026-08-29:} on the card's polish-pass audit (repair P1-R46,
finding F11, applied to \texttt{research/model\_v4/proofs/P1\_proof.md} Step 18) the Kakutani
conclusion recorded above is withdrawn as unestablished --- convex values and a closed graph of
$\mathfrak T$ would need a lemma the proof does not supply, so no Kakutani conclusion is drawn
there. The step stays a remark outside the claim, no step reads it, and nothing the proposition
rests on changes.

\medskip
\noindent\emph{Part F: Assumption~\ref{asm:A8} and both cells on path.}

\smallskip
\noindent\textbf{Step 19 (at an equilibrium at which Assumption~\ref{asm:A8} holds, both cells carry
positive probability).}
At $k^\star$, hypothesis (P-7) makes $\mathcal C_F = \{D=1\}$ and $\mathcal C_P = \{D=0\}$ exclusive
and exhaustive, so $\Prb(\mathcal C_F) = \Omega(\kappa,\tau,T)$ and
$\Prb(\mathcal C_P) = 1-\Omega(\kappa,\tau,T)$, with $\Omega$ evaluated under the equilibrium plan
map $j^\star(\cdot\,;k^\star)$ of Step~17(i). The value of $\Omega$ is unaffected by the choice
between that map and the induced $j_{k^\star}$, since the two can differ only at the finitely many
cutoff points, a $\Phi_s$-null set. Assumption~\ref{asm:A8} asserts $0 < \Omega < 1$, so both
probabilities are strictly positive: both cells are reached with positive probability under the
equilibrium, that is, both are on path. This is also the condition under which the cell averages
$M_F$ and $M_P$ of Definition~\ref{def:premium} are defined, and clause (TR-iii)'s implication
$D = 1 \Rightarrow a = 1$ makes the flagged cell an engagement cell throughout.

\smallskip
\noindent\textbf{Step 20 (what Assumption~\ref{asm:A8} does and does not do, and the single-threshold
restatement).}
Assumption~\ref{asm:A8} is a restriction on an \emph{equilibrium object}: $\Omega$ is computed at
$k^\star$, not from primitives. No step above rules out $\Omega(k^\star) \in \{0,1\}$, so this
proposition does not produce an equilibrium satisfying that assumption; it states that if the
constructed equilibrium satisfies it then both cells are on path, which is why the addendum is
written ``at any such equilibrium at which Assumption~\ref{asm:A8} holds''. Read literally, Step~19
is close to a restatement of the assumption, and its content is the consistency check that the
partition of Lemma~\ref{lem:D1} is non-degenerate at the fixed point.

The restatement that gives the assumption something to bite on is the following. Suppose in addition
(a) the upper-set engagement-flag hypothesis --- the flags $a_j$ equal $1$ exactly on an upper set of
the ordered menu, so that there is $j_a$ with $a_j = 1$ for $j \ge j_a$ and $a_j = 0$ for $j < j_a$;
(b) $\partial_s B_j \ge 0$ on Voice plans, which is clause (TR-iii) and hence already carried by
hypothesis (P-13); and (c) H-ord, Voice stake monotonicity across plans: for Voice plans $j' > j$,
$B_{j'}(s,d) \ge B_j(s,d)$ at every $(s,d)$. Of these only (b) is supplied by the standing
hypotheses; (a) and (c) are the two additions the proposition's final paragraph names, and this is
the one step that consumes them. Definition~\ref{def:plans} orders the menu by aggressiveness without
tying that order either to the engagement flags or to the stake path, which is why neither follows.
Then the flagged set
\[
\bigl\{s : a_{j^\star(s;k^\star)} = 1 \ \text{ and } \ B_{j^\star(s;k^\star)}(s,H-T) \ge \tau\bigr\}
\]
--- the equivalence $f_j \le H \iff B_j(s,H-T) \ge \tau$ being part~(b) of Lemma~\ref{lem:D1} --- is
an upper interval of signals. The first condition defines an upper set because
$j^\star(\cdot\,;k^\star)$ is weakly increasing, by construction in \eqref{eq:P1-jstar}, together with
(a); and within that set the map $s \mapsto B_{j^\star(s;k^\star)}(s,H-T)$ is weakly increasing,
because it increases in the signal at a fixed plan by (b) and increases across plans by (c). Writing
$s_F(k^\star)$ for the infimum of that upper interval,
\begin{equation}
\label{eq:P1-Omega}
\Omega \;=\; 1 - \Phi_s\bigl(s_F(k^\star)\bigr),
\end{equation}
and Assumption~\ref{asm:A8} is equivalent to $s_F(k^\star)$ being finite and strictly above
$-\infty$. Without (a) and (c) the flagged set need not be an interval and $\Omega$ need not be a
single-threshold object at all.

\medskip
\noindent\textbf{Scope of the conclusion.}
Six disclaimers travel with the construction above and are stated here rather than left to be
inferred.

First, \emph{uniqueness is not claimed}, in any form: not uniqueness of $k^\star$, not local
uniqueness, not uniqueness of the induced price system, and not uniqueness within a collapse face.
Brouwer's theorem is an existence theorem, and no step bounds $\lVert D_k\Tmap\rVert$ ---
Assumption~\ref{asm:AGE} (AGE) does that, on a region, and only Proposition~\ref{prop:C1} consumes it.

Second, the addendum is conditional at the equilibrium and nowhere else.
Assumption~\ref{asm:A8} is imposed at the fixed point, Step~20 shows that no step produces a fixed
point with $\Omega \in (0,1)$, and the existence half of the proposition consumes that assumption
nowhere. That an equilibrium satisfying it exists at any given parameter vector is not claimed.

Third, the boundary values of $\kappa$ are handled by extension and not by restriction. At
$\kappa \in \{0,1\}$ a pooled history needing a mark outside the surviving noise support is null
under every plan profile, so it is off nature's path rather than off the players', and
Definition~\ref{def:cpbe}\,(vi) asks nothing of it; the earlier claim that a limiting belief exists at
\emph{every} pooled history is false at $\kappa = 1$ and is withdrawn. No positivity of the flagged
cell at $\kappa = 1$ is asserted anywhere above, and none follows.

Fourth, the equilibrium is an object built from one fixed convention. The reference-belief root of
Step~9(c) can change the payoff on a $\Phi_s$-null set of signals and so can move a component of
$\Tmap$; a different admissible reference belief yields another equilibrium rather than a failure of
this one, and the proposition does not claim convention-independence. In the same spirit, neither the
date at which the engagement cost is incurred nor the reading of \eqref{eq:Y}'s takeover-branch price
is settled here: Remark~\ref{rem:Ctiming} records the first as open between two readings and
hypothesis (P-10) is what makes the conclusion the same under both, while Step~5 records the second
and shows its conclusion survives either.

Fifth, the flagged order is optimal but not \emph{uniquely} optimal. Step~12 proves the opposite of
uniqueness: on each deviation class the flagged price is invariant and the order cancels, so every
element of $\mathcal Q_j(s)$ delivers the same continuation and the blockholder is exactly
indifferent over the whole action set. The specified $Q_j^F(s)$ is \emph{a} maximiser, which is all
requirement~(ii) of Definition~\ref{def:cpbe} asks; the model does not pin the flagged order by
incentives. This is the round-2 face of the full separation Assumption~\ref{asm:A7J} (A7-J) buys: once the
filing reveals the signal and the price is competitive, no informational rent survives to make one
order strictly better than another.

Sixth, two of the hypotheses are sufficient rather than necessary and are not verified menu by menu.
Assumption~\ref{asm:A7J} (A7-J) is exhibited as satisfiable --- the pro-rata single-Voice menu with terminal
target strictly increasing on the whole line satisfies it, and it is the paper's pinned instance ---
but nothing above establishes it on a general finite menu, and the discussion after that assumption
records its failure boundary. Hypothesis (P-9) is \emph{a} sufficient condition delivering
requirement~(ii) at flagged nodes and is not claimed to be the weakest, the off-path-belief route
permitted by Definition~\ref{def:cpbe}\,(vi) being untouched by the plan-only perturbation used here;
hypothesis (P-10) inherits the same three disclaimers, being vacuous on any single-Voice menu and a
genuine restriction only on menus whose Voice plans share a pooled path, where no step above verifies
it. Finally, Assumption~\ref{asm:A6} is not claimed to be derivable at this level of generality: its
ordering and self-map content is derived in Step~13, while the bracket of Step~14 and the continuity
of Step~15 are assumed, and Remark~\ref{rem:A6record} carries the measured record that the continuity
clause fails at the implemented calibration for the $\Theta$ this argument constructs.

\medskip
With those six scope statements attached, the proof of the proposition is complete.
\end{proof}
